\section{Core reduction searches and a bootstrap obstruction}

No binary blank-tape machine below 245 working states was certified in this
component of the investigation. The new contribution is a collection of
broader finite reduction searches and a rigorously checked example showing
that even universally correct register operations need not preserve blank
initialization.

\subsection{The spare-zero resource variant}

Write
\[
 B(v)=0\,1^{v+1},\qquad C(v)=1\,0\,1^v,
\]
and extend both encodings by concatenation to finite lists of natural
numbers. The elementary identity
\[
 C(A)1=1 B(A)
\]
follows by induction on the list $A$. The routed machine uses this identity
when its carry scan temporarily converts blocks of type $B$ into type $C$.
The ordinary return scan and rollback scan restore their respective frames.

Consider the exact old-label 245-state table, and change only
\[
 255:\quad 0\mapsto 1 L 255
 \qquad\hbox{to}\qquad
 255:\quad 0\mapsto 0 L 255.
\]
All other rows, including the start state 95, are initially retained. This
modification will be called the spare-zero table; it is not a successful
consistency construction.

At the right end of an increment, the relevant local rows are
\[
\begin{array}{c|cc}
 &0&1\\\hline
 75&1 R 255&1 L 63\\
 255&0 L 255&1 L 131\\
 131&0 L 48&0 L 112.
\end{array}
\]
The sole changed write leaves a farther blank cell zero instead of making
it one. Thus the new increment appends $B(0)$ where the original appends
$B(1)$. The zero-decrement case reaches state 255 on the first one-cell of
the following nonempty block and therefore does not use the changed row.
Its rollback remains unchanged. Positive decrement is likewise unchanged.

The universal local proof was replayed on the actual modified table. All
four one-block inductions, the list-word identity, all three complete
increment/decrement macros, and marking contracts pass. The only adaptation
to the macro proof driver replaces the increment target's appended $B(1)$
by $B(0)$. No execution rule is weakened. An independent literal check also
passes all 8385 archived concrete instruction cases after changing exactly
2599 increment targets in this way.

This resource flexibility is legitimate inside the runtime invariant:
spare blocks may have arbitrary natural values, and an increment still
preserves a nonempty finite spare suffix. A complete whole-frame symbolic
replay on the actual modified table also passes all 8391 guarded cases,
using 2312645 symbolic rules and covering all 490 transition rows. It
quantifies independent arbitrary primary register values, arbitrary gap
$64+n$, and an arbitrary nonempty spare suffix, with minimum progress 13.
The new proof is an auditable mathematical symbolic check, not a new
proof-assistant endpoint. These results do not justify the blank
initialization of the modified machine.

\subsection{An explicit infinite continuation after 78 steps}

\paragraph{Blank-ray lemma.}
Suppose a binary Turing machine is in state $q$ at position $h$, every cell
at or to the right of $h$ is zero, and
\[
 \delta(q,0)=(w,R,q)
\]
for $w\in\{0,1\}$. Then it does not halt. More precisely, after $n$ steps
its state is $q$, its head is $h+n$, and it has changed exactly the cells
$h,\ldots,h+n-1$ to $w$; all cells farther right remain zero. The analogous
leftward statement also holds.

\paragraph{Proof.}
At $n=0$ the description is the premise. Assuming it at $n$, the scanned
cell $h+n$ is zero, so the stated transition writes $w$, moves to $h+n+1$,
and keeps state $q$. This proves the description at $n+1$. In particular,
every finite prefix extends by one further nonhalting transition.

Two independent exact interpreters, using sparse sets and dense arrays,
check that the spare-zero table started at 95 on blank tape reaches after
exactly 78 transitions
\[
 q=42,\qquad h=4,\qquad
 \{x:T(x)=1\}=\{-3,-2,2\}.
\]
The literal row $\delta(42,0)=(1,R,42)$ and the displayed complete tape
support satisfy the blank-ray lemma. Consequently, after an additional
$n$ steps,
\[
 q=42,\qquad h=4+n,\qquad
 \{x:T(x)=1\}=\{-3,-2,2\}\cup\{4,5,\ldots,3+n\}.
\]
This is an elementary proof of nonhalting, not a guess based on a long
simulation or an unverified absence of a supported root.

The same certificate generator checks every one of the 245 working states
as a possible blank start. Exactly 243 starts reach a blank-ray self-loop.
The other two starts halt:73 after 5380 transitions and 83 after 5379. There
are no unresolved starts. Every endpoint is independently reproduced with
the dense interpreter, and each escape endpoint is checked against its
actual table row and entire finite one-cell support. Thus changing only the
start state does not repair the intended calibration of this particular
spare-zero table.

The distinction between truth and provability matters here. If ZFC is
consistent, the external statements ``this escaping machine does not halt''
and $\operatorname{Con}(\mathrm{ZFC})$ are both true. The failure concerns
the requested provable equivalence: under the usual hypotheses of the
second incompleteness theorem, a consistent ZFC cannot prove its consistency
by combining a proved nonhalting statement with such an equivalence. No
absolute consistency or inconsistency assertion is made.

\subsection{Broader finite reduction families}

The single-incoming-edge search permits deletion of any state other than
start 95 or dispatcher 96 having exactly one nonself incoming edge. That edge
may choose an arbitrary write bit, either direction, and any retained working
state as successor. All 137 eligible states were screened, giving
\[
 137\cdot 244\cdot 4=133712
\]
concrete rewrite candidates. None survives the 8385 instruction fixtures.
A second pass allows translated dispatcher frames, a different safe gap,
and arbitrary nonempty spare contents; it also has no survivor.

The contextual family allows 2--6 incoming edges and permits a different
replacement for each caller. It first retains individually successful edge
edits, then tests every combination when all incoming edges have an
alternative. The run completes 161040 individual screens and 1908120
composed rewrites, again without a survivor. This is a restricted
compositional family: an edit requiring simultaneous mutually repairing
changes may be missed by its individual-edge prefilter.

A separate synthesizer replaces the four end-handler roles 75,112,131,255
by three unrestricted two-row states. With five successor choices, the
whole-instruction check explores 595777 partial-table nodes. A broader
seventeen-successor family explores 7534057 nodes. Both have no survivor.
The endpoint is the complete instruction frame, so intermediate tape edits
and head positions may differ from the inherited implementation.

Another synthesis removes state 63 and replaces the return roles 48,112,131
by fresh rows, while retargeting 75/read 1. The all-left family explores 231420
partial-table nodes; allowing both head directions explores 8312852. Both
finish without a survivor. The spare-zero targets were also used in separate
single-edge, contextual and end-handler searches; none yields a 244-state
candidate.

Every search has explicit finite bounds. Single-edge and contextual quick
screens use 6000 transitions and final affected-case screens 10000. Whole
instruction synthesis uses 4000. Concrete tape coordinates range from$-64$
to 1983. The local end-handler screen uses 56 finite contexts and 120 steps,
allowing an appended spare ofvalue 0 through 4. The archived baseline itself
passes all 8385 fixtures in at most 2686 steps.

These outcomes are not lower bounds on consistency-machine size. A correct
candidate might require more time on a test input, different representation
invariants, a new source program, a successor outside the chosen pool, or a
synergistic edit excluded by the contextual prefilter. The results provide
reproducible search evidence and a concrete bootstrap warning; they do not
establish minimality.
